problem:  
Let \( G \) be a compact, connected, semisimple Lie group with Lie algebra \( \mathfrak{g} \), and let \( H \subset G \) be a closed subgroup with reductive decomposition \( \mathfrak{g} = \mathfrak{h} \oplus \mathfrak{m} \).  
Fix a left-invariant, bracket-generating distribution \( \Delta \subset TG \) such that \( \Delta_e = \mathfrak{m}' \subset \mathfrak{m} \), equipped with a sub-Riemannian metric induced by the restriction of the negative Killing form.

Consider the normal Hamiltonian 
\[
H(p_g) = \tfrac{1}{2}\langle p_g , \Pi_{\mathfrak{m}'}(p_g)\rangle_{\mathfrak{g}^*},
\]
defined on \( T^*G \), where \( \Pi_{\mathfrak{m}'}: \mathfrak{g}^* \to (\mathfrak{m}')^* \) is the orthogonal projection with respect to the Killing form pairing on \( \mathfrak{g} \simeq \mathfrak{g}^* \).  
The left \(G\)-action on \(T^*G\) is Hamiltonian with moment map \( \mu_L(g,p_g)=\operatorname{Ad}_g^* p_g \).

Let \( \mathcal{C}\subset T^*G \) be the coadjoint-invariant subset
\[
\mathcal{C} = \{(g,p_g)\in T^*G \mid p_g \in (\mathfrak{m}')^* \}.
\]
The normal sub-Riemannian geodesic flow preserves \(\mathcal{C}\).

**Open problem:**  
Assume that the sub-Riemannian Hamiltonian \(H\) is **Liouville integrable** on each connected component of \(\mathcal{C}\). Is it necessarily true that such integrability derives from a **Poisson-commutative algebra of \(G\)-invariant functions on \(\mathfrak{g}^*\)**, restricted to \((\mathfrak{m}')^*\)?  
Equivalently:  
> Must every Liouville-integrable left-invariant sub-Riemannian geodesic flow on a compact semisimple Lie group arise via restriction of a Mishchenko–Fomenko-type integrable structure on \( \mathfrak{g}^* \)?  
If not, can one construct an explicit example where integrability holds but cannot be obtained from such a restriction, indicating a genuinely non-symmetric (purely dynamical) integrable mechanism?

---

Why is it a "good" problem:  
This reformulation precisely targets **the structural source of integrability** in left-invariant sub-Riemannian systems, stripped of ambiguous geometric terms. It now asks whether **all Liouville-integrable sub-Riemannian flows arise from algebraic Poisson-commutative families on \(\mathfrak{g}^*\)**, or whether fundamentally new integrable behaviors exist that violate symmetry-based expectations.

It is "good" because:

* **Foundational clarity:** The problem sits at the intersection of sub-Riemannian geometry, Poisson geometry, and the theory of integrable systems, framed in rigorous algebraic terms (Poisson-commutative families, moment maps, Mishchenko–Fomenko construction).
* **Deep classification potential:** A positive answer would imply that every known integrable invariant sub-Riemannian flow stems from hidden group symmetries; a counterexample would signal an entirely new nonholonomic integrability mechanism.
* **Mathematical precision with broad reach:** The question demands sharp control-theoretic and algebraic tools, inviting methods from representation theory, symplectic reduction, and geometric control.
* **Conceptually simple yet profound:** In a single line—“Does integrability always come from symmetry?”—it encapsulates a bridge between analytic and algebraic structures in modern geometry.

This problem is therefore mathematically valid, research-level, structurally clear, and deeply consequential—a paradigmatic example of “narrow entrance, profound depth.”